% !PRINTED 573
\weberpage{573}
%<<CONT
and $\tau_f$ is congruent with $\tau$. Accordingly we obtain the
congruence:
\[
[\Pi(t, t_1, t_2, \ldots, t_n)]^p \equiv \Pi(t^p, t_1^p, t_2^p, \ldots, t_n^p)\ (\mathrm{mod}\ \mathfrak p).
\]
By \S.\,150, 4. the theorem is thus proved:

2. \gesperrt{The form}
\[
\Pi(t) = (t - \tau)(t - \tau_1) \ldots (t - \tau_{f-1})
\]
\gesperrt{is congruent, with respect to the modulus $\mathfrak p$, to
an integral and homogeneous form of degree $f$ in $R$ in the variables
$t, t_1, t_2, \ldots, t_n$.}

This integral rational form, whose coefficients are completely
determined up to multiples of the prime $p$, we shall denote by $P(t)$.

Then from (13) follows
\weq{14}{\Phi(t) \equiv P(t)\, \Phi_0(t)\ (\mathrm{mod}\ \mathfrak p),}
and since now $\Phi(t)$ is an integral function in $R$, raising to the
power $p$ yields:
\[
\begin{aligned}
\Phi(t^p, u^p) &\equiv P(t^p, u^p)\ \Phi_0(t^p, u^p) \\
[\Phi(t, u)]^p &\equiv [P(t, u)]^p\ [\Phi_0(t, u)]^p
\end{aligned}
\qquad (\mathrm{mod}\ \mathfrak p).
\]
But because the left-hand sides are congruent with respect to the
modulus $\mathfrak p$, it follows that
\[
[P(t)]^p\ [\Phi_0(t, u)^p - \Phi_0(t^p, u^p)] \equiv 0\ (\mathrm{mod}\ \mathfrak p),
\]
and from this, since $P(t)$ is not divisible by $\mathfrak p$ (because
the coefficient of $t^f$ has the value 1):
\weq{15}{[\Phi_0(t, u)]^p \equiv \Phi_0(t^p, u^p),}
from which it follows by \S.\,150, 4. that $\Phi_0(t)$ too is
congruent to an integral rational form with respect to the modulus
$\mathfrak p$.

Accordingly we may in (14) also take $\Phi_0(t)$ to be an integral
form in $R$, and then by \S.\,140, 3. the congruence (14) must hold
not only for the modulus $\mathfrak p$, but for the modulus $p$.
From this we obtain the theorem:

3. \gesperrt{Among the forms $\Phi(t)$ which for $t = \tau$ pass over
into a functional divisible by $\mathfrak p$, the form $P(t)$, of
degree $f$ with respect to $t$, is the form of lowest degree, and if
$\Phi(t)$ is any one among them, then an integral function $\Phi_0(t)$
in $R$ can be so determined that}
\weq{15}{\Phi(t) \equiv P(t)\, \Phi_0(t)\ (\mathrm{mod}\ p)} % [sic] number (15) reused for a second equation on this page, as printed
\gesperrt{holds.}
